% !TeX root = ../papers/nc-p3.tex
\section{階数一特異点の局所長と固定像直線配置の排除}
\label{sec:cofactor-singular-curve}
定理\ref{thm:generically-reduced-rank-two-curves}の仮定と記号を用いる。
正規化 $\nu:\widetilde C\to C$ 上の $\nu^*W$/捩れの総次数を $\beta$、
$X$ のうち $C$ 上の印数を $m$、$C$ 上の相異なる階数一点の総数を $s$ とする。
一般階数は各成分で二、十成分に共通零点はないので $s$ は有限である。
\begin{proposition}\label{prop:cofactor-singular-length}
次の二つの上界を得る。
\[
 6d+2\chi(\mathcal O_C)+\beta+2(s-m)\le 2m,\qquad
 \beta+2(s-m)\le 2h^1(\mathcal O_C(3))-2h^1(\mathcal I_C(3)).
\]
\end{proposition}
\begin{proof}
$D=\ker(\mathcal O_Z\to\mathcal O_C)$ の局所長を $a_p$ と書く。
前節の $0\to T\to K\to\mathcal I_Z\to0$ において、純一次元 $T$ は余次元二なので
$\mathcal Ext^1(T,\omega)=0$。従って $(\bigwedge^2 M)(6)\to W(6)$ の核は
$\mathcal Ext^2(\mathcal I_Z,\omega)=\mathcal Ext^3(\mathcal O_Z,\omega)
=\mathcal Ext^3(D,\omega)$ であり、各点で長さ $a_p$ を持つ。
最後の同型は被約純曲線がCohen--Macaulayであることによる。
階数一点で $\bigwedge^2 M$ のファイバーは三次元だから
$a_p\ge3-\mu_p(W)$。ここで $\mu_p$ は局所最小生成元数である。

$W$ の一般大域切断を選び、正規化の線束への像が全ての特異点・非自由点・階数一点の上で
零でないようにする。これは有限個の一次条件の回避である。
$\mathcal O_C\hookrightarrow W$ の余核長 $b$ は滑らかな零点で総数 $\beta$ を与え、
その他に各階数一点で少なくとも $\mu_p(W)-1$ を与える。
特にそこでは $a_p$ とこの追加分の和が二以上になる。
$C$ の外の印は二重点の長さ四を与えるから
$a+b\ge 4(20-m)+2s+\beta$。八十のEuler式に代入して第一式を得る。
三次式はこれら $m$ 点へ全射するので
$m\le h^0(\mathcal O_C(3))-h^1(\mathcal I_C(3))$、従って第二式が従う。
\end{proof}

\begin{theorem}\label{thm:cofactor-singular-special}
一般被約な階数二曲線全体で一般像二平面が各成分ごとに一定となることはない。
特に必ず $h^1(\mathcal O_C(3))>h^1(\mathcal I_C(3))$ である。
\end{theorem}
\begin{proof}[直線と多重点の方向]
固定像なら各成分の二次Veronese像の線形包はベクトル次元三以下なので成分は直線である。
実際非直線曲線の一般三点は非共線で、さらに別の点の平方はその三平方の包に入らない。
$N(x)=P(x^2)$ と書き、十次元の同型 $P$ を用いると、各入力直線の二次平方の包は
出力直線の二次平方の包へ同型に移る。異なる直線は異なる直線に対応する。
二つの平方の包の共通部分は元の二空間の共通部分の対称平方なので、交点は階数一で、
交点と非交点の区別が保たれる。
三本が同一平面内で三角形なら出力でも三角形をなし、三本が一点に集まるときは出力でも一点に集まる。
後者の場合、三本の平方の包の和の次元は共面なら六、非共面なら七である
（共通点の平方、その点と三方向との積、三方向の平方に分ける）。
いずれも共面な三直線があれば一平面の全平方が一平面の全平方に移る。
すると入力平面で固定核列が消え、定理\ref{thm:opposite-column-divisor}に反する。
従って入力・出力とも三本の共面直線はない。

一点 $p$ に $k\ge4$ 本が集まるとする。$p$ と出力点を第一座標、入力を $(t,u)$, $u\in k^3$
と書く（本数の $k$ と係数体は文脈で区別する）。方向をそれぞれ $d_i,e_i\in\mathbb P^2$ とする。
どちらでも三方向が共線になることはない。下三行三列の $tu$ 係数を $H(u)$ とすると
$H(d_i)\in\langle e_i^2\rangle$。任意の四方向の線形関係は四係数とも非零であり、
出力の四平方は独立である。実際各一点を避け他三点を覆う二直線の積を取れる。
従って $H=0$。$P$ は可逆なので接空間の間に誘導する三次元線形写像は可逆、
さらに商の六次元写像も可逆である。基底を合わせて行列を
\[
 \begin{pmatrix} * & (tu+B(u))^t\\tu+B(u)& F(u^2)\end{pmatrix}
\]
とできる。ただし $B$ は三個の二次式、$F$ は可逆である。各直線について
$e_i=d_i$, $B(d_i)\parallel d_i$, $F(d_i^2)\parallel d_i^2$ が成り立つ。
$F$ がスカラーなら全印で出力の下三座標は $u$ と比例し、
独立な混合積の間に交代関係を生じる（$u=0$ でも関係は消える）。
従って $F$ は非スカラーである。

これらの方向のうち六点が同じ二次曲線に載ることはない。実際三共線がないのでその二次曲線
$E$ は滑らかであり、任意の五平方は独立（二本の直線で他の四点を覆う）。
六平方は五次元を張り、その唯一の関係は全係数が非零。
従って $F$ はこの五次元上スカラーで、$E$ 全体で $F(d^2)\parallel d^2$。
$B(d)\bmod d$ は $E\simeq\mathbb P^1$ 上の束
$(\mathcal O^3/\mathcal O(-2))(4)\simeq\mathcal O(5)^2$ の切断である。
この分解は $(s^2,st,t^2)$ の二つの一次関係で従う。六零点があれば恒等的に零で、
$E$ 上の全方向の直線で $N$ の像は二次元以下、階数二曲面排除に反する。
七方向があれば最初の六平方は従って独立、これを固有基底とする。
第七平方の展開は少なくとも六係数が非零である。さもなければ六点以下に線形従属が生じる
（六平方の従属は六点が二次曲線に載ることと同値）。すると六固有値は同じで矛盾する。
よって $k\le6$ で、六方向のとき二次評価は全射である。
\end{proof}

\begin{proof}[局所長と全直線の印数]
全直線数を $r$ とする。各交点 $p$ における $k$ 本の方向の $j$ 次評価階数を $h(j)$ とする。
曲線は交点近傍で方向点のアフィン錐だから、正規化との局所余長は
$\delta=\sum_{j\ge0}(k-h(j))$ である。
各枝の $W$ の線束は自明で、その六生成元の枝での値は出力直線のPlücker列に比例する
（核と像の外積は一定の双対で交換できる）。交点では各列は共通点と方向 $e_i$ の外積となる。
上で $k\ge4$ では $e_i=d_i$ とでき、$k=2,3$ でも双方の方向が独立なので同じことができる。
各枝の定数倍を選べば、局所完備化での $W$ は正規化の中で方向の一次評価列が生成する加群である。
各次数 $j$ の係数の階数は $h(j+1)$ なので、$W$ の正規化との余長は
$\delta_W=\delta-(k-1)$。ゆえに八十のEuler式は
\[
      a=80-8r+\sum_p(2\delta_p-k_p+1)
\]
となる（和は相異なる交点）。

$k=2,3$ なら $\delta=k-1$, $\mu_p(W)=k$ より $a_p\ge3-k$。
$k\ge4$ では上記行列の単元第一行列要素を消すと残余三行の各成分は二次以上で、
その二行小行列式は四次以上である。従って $Z$ は $p$ の四乗肥厚を含み、
$a_p\ge20-\sum_{j=0}^3h(j)$。$k=4,5,6$ では $h(0)=1,h(1)=3,h(j)=k$ ($j\ge2$)。
以上の全場合で
\[
                      4-a_p+2\delta_p-k+1\le2k .
\]
実際 $k=4,5,6$ の $a_p$ の下界は順に $8,6,4$ である。
滑らかな印で $a_p\ge2$、曲線外の印で $a_p\ge4$。
各直線の印数の総和を $M$、印でない交点の集合を $I$ とすれば、これらをEuler式に入れて
$8r\le2M+\sum_{p\in I}(2k_p-4)$。
各直線上で固定像二行の行列式は非零四次式なので階数一点は高々四点、
従って $M+\sum_{p\in I}k_p\le4r$。二式より $I$ は空で $M\ge4r$。
一方、印は各直線に高々三つであり矛盾する。
これで各成分の像が一定の場合を尽くした。従って $\beta>0$ であり、
命題\ref{prop:cofactor-singular-length}から最後の主張も従う。
\end{proof}
